\section{Periods and Riemann Bilinear Relations}
\label{sec:periods}

For a singular one-cycle \(\gamma\) and \(\omega\in\HX\), define
\[
  \Pi_X(\gamma)(\omega)=\int_{\gamma}\omega.
\]
Since holomorphic one-forms are closed, the integral only depends on the
homology class of \(\gamma\). This gives the period pairing of
Definition~\ref{def:period-pairing}.

\begin{lemma}[Homology invariance of periods]
\label{lem:period-homology-invariance}
\notready
If \(\gamma\) and \(\gamma'\) represent the same class in \(H_1(X,\Z)\), then
\[
  \int_{\gamma}\omega=\int_{\gamma'}\omega
\]
for every \(\omega\in\HX\).
\lean{JacobianChallenge.Periods.period_homology_invariance_statement}
\end{lemma}

\begin{proof}
The difference \(\gamma-\gamma'\) is a boundary. Holomorphic one-forms are
closed. Stokes' theorem gives that the integral of a closed form over a
boundary vanishes.
\end{proof}

\begin{classicalinput}[Riemann bilinear relations]
\label{input:riemann-bilinear}
\notready
Let \(a_1,\dots,a_g,b_1,\dots,b_g\) be a symplectic basis of
\(H_1(X,\Z)\). For holomorphic one-forms \(\omega,\eta\),
\[
  \int_X \omega\wedge\eta
  =
  \sum_{i=1}^{g}
  \left(
    \int_{a_i}\omega\int_{b_i}\eta
    -
    \int_{b_i}\omega\int_{a_i}\eta
  \right).
\]
In particular, the period vectors of a basis of \(\HX\) have full real rank
\(2g\).
\lean{JacobianChallenge.Periods.period_vectors_linearIndependent_of_symplectic}
\end{classicalinput}

\begin{proofsketch}
Cut \(X\) along the symplectic basis to obtain a polygon. On the polygon every
holomorphic one-form has a primitive. Apply Stokes' theorem to the boundary
integral and identify opposite sides. Taking \(\eta=\overline{\omega}\) in
the Hermitian variant gives positivity:
\[
  i\int_X \omega\wedge\overline{\omega}>0
\]
for \(\omega\ne 0\). This positivity forces nondegeneracy of the period
matrix and full rank of the period lattice.
\end{proofsketch}

\begin{proof}[Proof of Theorem~\ref{thm:period-lattice}]
Choose a \(\C\)-basis \(\omega_1,\dots,\omega_g\) of \(\HX\), so
\(\HdualX\cong\C^g\). The images of the \(2g\) symplectic cycles under
\(\Pi_X\) are the columns of the period matrix. By the Riemann bilinear
relations, these \(2g\) vectors are linearly independent over \(\R\). Since
\(\HdualX\) has real dimension \(2g\), their \(\Z\)-span is a full lattice.
\end{proof}

\begin{formalizationnote}
The current project has basis-aligned period subgroup files. The deepest
remaining content is the full-rank statement, routed through declarations such
as \code{periodSubgroup_isZLattice},
\code{periodSubgroup_spans_real}, and
\code{period_vectors_linearIndependent_of_symplectic}.
\end{formalizationnote}
